\section{Related Work}
\label{sec:related-work}

\noindent\textbf{eBPF optimization.} K2 showed that BPF bytecode has substantial optimization headroom before load~\cite{xu2021synthesizing}. Merlin extended this view across LLVM IR and BPF bytecode, emphasizing that eBPF optimization is naturally multi-tier~\cite{mao2024merlin}. \sys differs in its evaluation target: it benchmarks agents that choose optimization actions for existing programs under real verifier, JIT, and workload feedback. A K2- or Merlin-style optimizer could be one backend action in \sys, but the benchmark's unit is the closed-loop tuning task.

\noindent\textbf{Verifier and safety systems.} PREVAIL, verifier-validation work, and verifier-correctness studies focus on proving or auditing eBPF safety~\cite{gershuni2019simple,sun2024validating,vishwanathan2023verifying}. Rex and KFlex change the programming or safety model to make kernel extensions safer and more usable~\cite{jia2025rex,dwivedi2024fast}. These systems motivate the need for strong acceptance oracles. \sys keeps the Linux verifier in the loop and treats verifier acceptance as one necessary component of optimization success.

\noindent\textbf{LLMs and eBPF.} Kgent and related systems ask whether language models can generate kernel extensions or eBPF programs~\cite{zheng2024kgent}. Verifier-repair benchmarks ask whether models can repair programs rejected by the verifier. \sys asks a different question: given an existing eBPF program that already belongs to an application, can an agent choose a performance-improving transformation without breaking verifier safety or workload behavior?

\noindent\textbf{Agent benchmarks.} Software-engineering and tool-use benchmarks such as SWE-bench and \(\tau\)-bench evaluate agents in executable environments with hidden tests, reliability metrics, or cost budgets~\cite{swebench2023,taubench2024}. Recent benchmark-audit work argues that incomplete graders, ambiguous specifications, and environment dependencies can change model rankings~\cite{swebenchsolved2025,benchmarkaudit2026}. \sys follows this line but changes the oracle. The benchmarked artifact is not a Python patch or a CUDA kernel alone; it is an eBPF optimization decision whose validity is decided by the Linux verifier, the kernel JIT, application startup, workload semantics, protected benchmark integrity checks, and post-hoc performance methodology.
